\documentclass[12pt, a4paper]{article}
\usepackage[utf8]{inputenc}
\usepackage[T1]{fontenc}
\usepackage[ngerman]{babel}
\usepackage{amsmath, amssymb, amsthm}
\usepackage{geometry}
\usepackage{graphicx}
\usepackage{xcolor}
\usepackage{tcolorbox}
\usepackage{enumitem}
\usepackage{listings}
\usepackage{fancyhdr}
\usepackage{titlesec}
\usepackage{mdframed}
\usepackage{tikz}
\usepackage{pgfplots}
\pgfplotsset{compat=1.18}

\geometry{a4paper, margin=2.5cm}

\definecolor{theoryblue}{RGB}{30, 90, 160}
\definecolor{exerciseorange}{RGB}{200, 80, 0}
\definecolor{solutiongreen}{RGB}{0, 120, 60}
\definecolor{hintgray}{RGB}{100, 100, 100}
\definecolor{codebg}{RGB}{245, 245, 245}

% Theory box
\tcbuselibrary{skins, breakable}
\newtcolorbox{theorybox}[1]{
  colback=theoryblue!5,
  colframe=theoryblue,
  fonttitle=\bfseries\color{white},
  title=#1,
  breakable
}

% Exercise box
\newtcolorbox{exercisebox}[1]{
  colback=exerciseorange!5,
  colframe=exerciseorange,
  fonttitle=\bfseries\color{white},
  title=Aufgabe #1,
  breakable
}

% Solution space box
\newtcolorbox{solutionspace}[1][6cm]{
  colback=white,
  colframe=black!20,
  fonttitle=\itshape\color{hintgray},
  title=Rechenraum,
  height=#1,
  breakable
}

% Hint box
\newtcolorbox{hintbox}{
  colback=solutiongreen!5,
  colframe=solutiongreen,
  fonttitle=\bfseries\color{white},
  title=Hinweis,
  breakable
}

% Julia code style
\lstdefinelanguage{Julia}{
  keywords={function, end, return, if, else, elseif, for, while, in, begin, let, do, try, catch, finally, using, import, export, mutable, struct, abstract, type},
  keywordstyle=\color{theoryblue}\bfseries,
  comment=[l]{\#},
  commentstyle=\color{hintgray}\itshape,
  stringstyle=\color{solutiongreen},
  morestring=[b]",
  sensitive=true
}
\lstset{
  language=Julia,
  backgroundcolor=\color{codebg},
  basicstyle=\ttfamily\small,
  breaklines=true,
  frame=single,
  rulecolor=\color{black!30},
  numbers=left,
  numberstyle=\tiny\color{hintgray}
}

\pagestyle{fancy}
\fancyhf{}
\rhead{\textcolor{hintgray}{\small Rocket Science Workbook}}
\lhead{\textcolor{hintgray}{\small \leftmark}}
\cfoot{\thepage}

\titleformat{\section}{\large\bfseries\color{theoryblue}}{}{0em}{}[\titlerule]
\titleformat{\subsection}{\normalsize\bfseries\color{exerciseorange}}{}{0em}{}


\begin{document}

\begin{center}
  {\LARGE \textbf{Kapitel 04: 2D-Physik der Rakete}}\\[0.5em]
  {\large \textcolor{hintgray}{Kräfte, Koordinatensystem, Zustandsvektor und Bewegungsgleichungen}}\\[0.3em]
  {\small \textcolor{hintgray}{Rocket Science Workbook — simulation-2D-jl}}
\end{center}
\vspace{1em}
\hrule
\vspace{1.5em}

% ─────────────────────────────────────────────
\section{Koordinatensystem}

\begin{theorybox}{2D Inertialsystem}
Wir verwenden ein \textbf{kartesisches Koordinatensystem} mit:
\begin{itemize}
  \item $x$: horizontale Distanz (Downrange), positiv = Richtung Zielorbit
  \item $y$: Höhe über dem Boden, positiv = nach oben
  \item $(0,\,0)$: Startrampe
\end{itemize}
Die Erde wird als flach angenommen (gültig für Reichweiten $\ll 6371\,\text{km}$).
\end{theorybox}

% ─────────────────────────────────────────────
\section{Zustandsvektor}

\begin{theorybox}{5-dimensionaler Zustand}
Der vollständige Zustand der Rakete zu einem Zeitpunkt $t$ ist:
\[
  \mathbf{s}(t) = \begin{pmatrix} x \\ y \\ v_x \\ v_y \\ m \end{pmatrix}
\]
\begin{itemize}
  \item $x, y$: Position (m)
  \item $v_x, v_y$: Geschwindigkeitskomponenten (m/s)
  \item $m$: aktuelle Gesamtmasse (kg)
\end{itemize}
\end{theorybox}

% ─────────────────────────────────────────────
\section{Kräfte}

\subsection*{Schwerkraft}
\[
  \mathbf{F}_g = \begin{pmatrix} 0 \\ -m\,g(y) \end{pmatrix}, \quad g(y) = g_0\left(\frac{R_E}{R_E + y}\right)^2
\]

\subsection*{Schubkraft}
Der Schubvektor zeigt in Richtung des Pitchwinkels $\theta$ (gemessen von der Horizontalen):
\[
  \mathbf{F}_T = T \begin{pmatrix} \cos\theta \\ \sin\theta \end{pmatrix}
\]
$\theta = 90°$ bedeutet senkrecht nach oben; $\theta = 0°$ bedeutet horizontal.

\subsection*{Luftwiderstand}
Der Drag wirkt \textbf{entgegen} dem Geschwindigkeitsvektor:
\[
  \mathbf{v} = \begin{pmatrix} v_x \\ v_y \end{pmatrix}, \quad
  |\mathbf{v}| = \sqrt{v_x^2 + v_y^2}, \quad
  \hat{\mathbf{v}} = \frac{\mathbf{v}}{|\mathbf{v}|}
\]
\[
  \mathbf{F}_D = -\frac{1}{2}\rho(y)\,|\mathbf{v}|^2\,C_d(Ma)\,A\;\hat{\mathbf{v}}
\]

% ─────────────────────────────────────────────
\section{Bewegungsgleichungen (ODE-System)}

\begin{theorybox}{Das vollständige DGL-System}
Newtons zweites Gesetz $\mathbf{F} = m\,\mathbf{a}$ liefert:
\[
\frac{d}{dt}\begin{pmatrix} x \\ y \\ v_x \\ v_y \\ m \end{pmatrix}
=
\begin{pmatrix}
  v_x \\
  v_y \\
  \dfrac{T\cos\theta - F_D\,\frac{v_x}{|\mathbf{v}|}}{m} \\[1em]
  \dfrac{T\sin\theta - F_D\,\frac{v_y}{|\mathbf{v}|} - m\,g(y)}{m} \\[1em]
  -\dot{m}
\end{pmatrix}
\]
\end{theorybox}

\begin{theorybox}{Sonderfall: $|\mathbf{v}| = 0$}
Wenn die Geschwindigkeit null ist, ist $\hat{\mathbf{v}}$ undefiniert. In diesem Fall ist $F_D = 0$. In der Implementierung: prüfe $|\mathbf{v}| < \varepsilon$ vor der Division.
\end{theorybox}

% ─────────────────────────────────────────────
\section{Rechenaufgaben}

\begin{exercisebox}{1 — Kräftezerlegung bei gegebenem Pitchwinkel}
Eine Rakete hat $T = 1\,400\,000\,\text{N}$, $\theta = 75°$, $m = 80\,000\,\text{kg}$.\\
$v_x = 100\,\text{m/s}$, $v_y = 800\,\text{m/s}$, $h = 5\,000\,\text{m}$.\\
$A = 10\,\text{m}^2$, $\rho(5000) \approx 0{,}7364\,\text{kg/m}^3$.

\textbf{a)} Berechne $|\mathbf{v}|$ und $Ma$:
\[
  |\mathbf{v}| = \underline{\hspace{3cm}}\,\text{m/s}, \quad Ma = \underline{\hspace{3cm}}
\]

\textbf{b)} Berechne $F_D$ (mit $C_d(Ma)$ aus Kapitel 03):
\[
  C_d = \underline{\hspace{2cm}}, \quad F_D = \underline{\hspace{4cm}}\,\text{N}
\]

\textbf{c)} Berechne die Schubkomponenten:
\[
  F_{T,x} = T\cos(75°) = \underline{\hspace{4cm}}\,\text{N}
\]
\[
  F_{T,y} = T\sin(75°) = \underline{\hspace{4cm}}\,\text{N}
\]

\textbf{d)} Berechne die Beschleunigungen $\dot{v}_x$ und $\dot{v}_y$:
\[
  \dot{v}_x = \frac{F_{T,x} - F_D\cdot\frac{v_x}{|\mathbf{v}|}}{m} = \underline{\hspace{4cm}}\,\text{m/s}^2
\]
\[
  \dot{v}_y = \frac{F_{T,y} - F_D\cdot\frac{v_y}{|\mathbf{v}|}}{m} - g = \underline{\hspace{4cm}}\,\text{m/s}^2
\]
\end{exercisebox}

\begin{solutionspace}[10cm]
\end{solutionspace}

\begin{exercisebox}{2 — Massenstrom und Schubberechnung}
Die Raketengleichung verbindet Schub $T$, Massenfluss $\dot{m}$ und Ausströmgeschwindigkeit $v_e$:
\[
  T = \dot{m}\,v_e
\]
Gegeben: $T = 1\,400\,000\,\text{N}$, $v_e = 3\,200\,\text{m/s}$ (typischer Wert für RP-1/LOX).\\[0.5em]
\textbf{a)} Berechne $\dot{m}$: \; $\dot{m} = \underline{\hspace{3cm}}\,\text{kg/s}$\\[0.5em]
\textbf{b)} Wie viel Treibstoff wird in 150 Sekunden verbraucht? \; $\Delta m = \underline{\hspace{3cm}}\,\text{kg}$\\[0.5em]
\textbf{c)} Eine Stage hat 30\,000 kg Treibstoff. Wie lange brennt sie?
\end{exercisebox}

\begin{solutionspace}[5cm]
\end{solutionspace}

\begin{exercisebox}{3 — Einen RK4-Schritt für das 2D-System}
Anfangszustand: $\mathbf{s}_0 = [0,\; 0,\; 0,\; 0,\; 100\,000]^\top$\\
Parameter: $T = 1\,400\,000\,\text{N}$, $\theta = 90°$, $\dot{m} = 437{,}5\,\text{kg/s}$, $A = 10\,\text{m}^2$, $h_\text{step} = 1\,\text{s}$\\[0.5em]
Vereinfachung: Nutze $\rho(0) = 1{,}225$, $C_d = 0{,}3$, $g = 9{,}81\,\text{m/s}^2$ (konstant).\\[0.5em]
\textbf{a)} Berechne $\mathbf{k}_1 = f(0, \mathbf{s}_0)$:
\[
  \mathbf{k}_1 = \begin{pmatrix} \dot{x} \\ \dot{y} \\ \dot{v}_x \\ \dot{v}_y \\ \dot{m} \end{pmatrix}
  = \begin{pmatrix} 0 \\ 0 \\ \underline{\hspace{2cm}} \\ \underline{\hspace{2cm}} \\ \underline{\hspace{2cm}} \end{pmatrix}
\]
\textit{Hinweis: Bei $|\mathbf{v}|=0$ ist $F_D = 0$.}\\[0.5em]
\textbf{b)} Berechne $\mathbf{s}^* = \mathbf{s}_0 + \frac{h}{2}\mathbf{k}_1$ (Zwischenpunkt für $k_2$).\\[0.5em]
\textbf{c)} Wie ändert sich $\dot{v}_y$ wenn $\theta = 45°$ statt $90°$ ist?
\end{exercisebox}

\begin{solutionspace}[10cm]
\end{solutionspace}

% ─────────────────────────────────────────────
\section{Verbindung zum Julia-Template}

\begin{theorybox}{Was du jetzt implementierst}
In \texttt{template.jl} (Kapitel 04) implementierst du:
\begin{enumerate}
  \item \texttt{gravity(y)} — höhenabhängige Schwerkraft
  \item \texttt{derivatives\_2d(t, s, params)} — die rechte Seite des DGL-Systems\\
        (gibt $[\dot{x}, \dot{y}, \dot{v}_x, \dot{v}_y, \dot{m}]$ zurück)
  \item Eine einfache Simulation: senkrechter Start, $\theta = 90°$, 1 Stufe
  \item Plot: Höhe, Geschwindigkeit, Beschleunigung über Zeit
\end{enumerate}
Die Funktion \texttt{derivatives\_2d} ist das Herzstück — sie wird von RK4 (Kapitel 02) aufgerufen.
\end{theorybox}

\end{document}
